\section{总结与延伸}

\subsection{核心结论}

\begin{enumerate}
\item 2025 年架构创新重新重要，是因为数据增长变慢、长上下文需求上升、推理成本变重。
\item Full Attention 表达力强，但长序列下计算和 KV cache 都贵。
\item Linear Attention 的核心是把平方注意力改写成接近线性的 recurrence 或状态更新。
\item Kimi Linear 的 KDA 用更细粒度 decay 改善有限 hidden state 的记忆利用。
\item Kimi 走 Hybrid Linear，DeepSeek 走 Sparse，MiniMax M2 回到 Full，说明 Attention 还未收敛。
\item FFN 已经被 MoE 改造，Attention 是 Transformer 里下一块可雕的区域。
\item 算法考古能把旧论文中被埋没的机制放到新硬件和新规模下重新验证。
\item 硬件亲和是架构研究底线：要能并行、能矩阵乘、能写 kernel、能 scale。
\item DeepSeek 的优势之一，是算法和 infra/硬件协同意识强。
\item 年轻研究者想做架构，需要接近 frontier lab 和算力环境，否则很难验证 scale。
\end{enumerate}

\subsection{开放问题}

最后保留开放问题，是因为 Attention 路线还没有收敛。Kimi、DeepSeek、MiniMax 的选择都不是终局，而是不同公司在不同约束下的阶段性押注。

\begin{itemize}
\item Kimi Linear 这类 Hybrid Attention 能否在更大规模持续接近 Full Attention 表现？
\item Sparse Attention 能否在保留长程检索能力的同时真正替代全局层？
\item 三比一的 hybrid 比例会成为长期经验，还是过渡方案？
\item Full Attention 是否会因为硬件继续优化而延长寿命？
\item MoE 之后，Attention 的突破会不会成为下一轮架构共识？
\item 架构研究如何在开源小规模和闭源大规模之间建立可靠验证桥梁？
\end{itemize}

\subsection{拓展阅读}

\begin{itemize}
\item Kimi Linear: An Expressive, Efficient Attention Architecture。
\item DeepSeek Sparse Attention / Native Sparse Attention 相关技术报告。
\item FlashAttention 系列：理解算法和硬件协同。
\item MoE、RoPE、NoPE、Mamba、DeltaNet、Gated Linear Attention 等架构材料。
\item Hazy Research 相关博客：高效序列建模和硬件友好架构。
\end{itemize}

\begin{importantbox}{最后的判断}
EP119 最值得留下的是一种研究姿势：不要只问“哪个 attention 更火”，而要问它解决哪个瓶颈、牺牲什么表达能力、如何在硬件上并行、能否通过 scaling ladder，以及是否能在真实 frontier 模型里稳定兑现。
\end{importantbox}

\end{document}
